% !TeX root = ../protokoll.tex
%  GRUNDLAGEN  (Abschnitt 1.2 der Anleitung: max. 2 Seiten!)
%  - nur SELBST erstellte Grafiken
%  - Formelsymbole müssen erklärt sein
%  - längere Herleitungen können in den Anhang
\section{Grundlagen}
\label{cap:Grundlagen}

\begin{comment}
    -2 Seiten lang (längere Herleitungen können in den Anhang gemacht werden)
    -nur SELBST erstellte Grafiken
    -Formelsymbole müssen erklärt sein
\end{comment}

Die theoretischen Grundlagen, die für das Verständnis der im Versuch untersuchten Vorgänge notwendig sind. Alle für die Auswertung benötigten Gleichungen müssen hier vorkommen.

\subsection{Grundgleichung}
\label{cap:Grundgleichung}
Beispielsweise gilt nach Newton
\begin{align}
    \vv{F} = m \cdot \vv{a}\;,
    \label{equ:Newton}
\end{align}
wobei $m$ die Masse, $\vv{a}$ die Beschleunigung eines Körpers und $\vv{F}$ die auf ihn wirkende Kraft beschreibt.

\subsection{Auswertungsformeln}
\label{cap:Auswertungsformeln}
Für den Zusammenhang zwischen den Messgrößen wird eine lineare Regression verwendet:
\begin{align}
    y = m \cdot x + b\;.
    \label{equ:Gerade}
\end{align}
Für Mittelwert, Standardabweichung und Fehler des Mittelwerts aus $N$ Messungen $x_i$ gilt:
\begin{align}
    \bar{x} &= \frac{1}{N} \sum_{i=1}^{N} x_i
    \label{equ:Mittelwert}\\
    \sigma &= \sqrt{\frac{1}{N-1} \sum_{i=1}^{N} \left(\bar{x} - x_i\right)^2}
    \label{equ:Standardabweichung}\\
    \Delta x &= \frac{\sigma}{\sqrt{N}}
    \label{equ:FehlerMittelwert}
\end{align}
Die Gauß'sche Fehlerfortpflanzung für $y = f(x_1, \dots, x_N)$ bei unkorrelierten Fehlern lautet:
\begin{align}
    \Delta y = \sqrt{\sum_{i=1}^{N}
    \left(\frac{\partial y}{\partial x_i}\right)^2 \cdot \left(\Delta x_i\right)^2}\;.
    \label{equ:Fehlerfortpflanzung}
\end{align}
